%!TeX root = ../incremental_SS_Translation.tex

\chapter{Śārīrasthāna 10:  On the analysis of the pregnant woman}


\section{Introduction}

\texttt{Draft text by Jan Gerris.
The Sanskrit text of this adhyāya is a lightly edited transcript of NAK 5-333,
waiting for full critical attention.}

\subsection{Literature}

Meulenbeld offered an annotated overview of this chapter and a
bibliography of earlier scholarship to 2002 and, in his notes, citations
of the parallel passages in the \CS.\fvolcite{IA}[259--261]{meul-hist}   

The first draft of the following translation was prepared by Jan
Gerris.
%, and later revised in collaboration Philipp Maas, John Nemec,
%Dominik Wujastyk, Madhusudan Rimal, Kenneth Zysk and Andrey Klebanov.

    
    
\section{Translation}

\begin{translation}
 
 
 \begin{tt}
     \raggedright
     \bigskip
     \marginpar{Draft tr.\ from here}
     
        
\item[1]

Now, we shall explain in detail the anatomy concerning the development
of the pregnant woman.

\item[2] \emph{vulgate only}

\item[3]

In that regard, beginning from the very first day of pregnancy, (she =
the expectant mother) should constantly remain joyful, clean, and
adorned in her attire, devoted to peaceful rites, auspiciousness, the
deities, Brahmins, and her mentors.

She should not touch those who are dirty, deformed, or have deficient
limbs.

Furthermore, she should keep at a distance those who smell foul or
appear repulsive.

She should avoid going outdoors, staying in empty houses, shrines,
cremation grounds, under sacred trees, as well as places that cause
anger and exhaustion, and activities such as speaking loudly and the
like.

She should not practice oil massages, body scrubs, and similar
treatments.

She should also not exhaust her body and must avoid the mentioned
previously activities .

She should arrange a bed and a sofa that are furnished with rich
coverings, not excessively high, equipped with a headrest, and are
sufficiently spacious.

She should prepare food that is agreeable to the heart, primarily liquid
and sweet, unctuous, prepared with digestive stimulants, and
well-cooked; this is the general rule until the time of delivery.

\item[4]

More specifically, during the first, second, and third months, the
pregnant woman should consume food that is primarily sweet, cooling, and
liquid.

In the fourth month, she should consume a diet of fresh butter derived
from milk.

In the fifth month, food mixed with milk and ghee (should be consumed).

In the sixth month, one should make her drink a rice gruel for half a
month, combined appropriately with ghee prepared with gokshura.

In the seventh month, the foetus is nourished in this manner by means of
substances prepared with Uraria picta and other medicinal plants.

In the eighth month, one should indeed administer a non-unctuous enema
to her using a decoction of jujube mixed with Sida cordifolia, Sida
rhombifolia, milk, curd, whey, oil, honey, ghee, and dill, in order to
clear out old fecal matter and to ensure the downward movement of the
vital breath. Thereafter, one should administer an unctuous enema with
milk and oil processed with sweet-tasting medicinal herbs.

Indeed, when the vital air is flowing in its natural downward direction,
she delivers smoothly.

And they remain free from any complications.

From this point forward, until the time of delivery, one should serve
her with unctuous gruels and the meat-juices of wild animals.

When treated in this manner, being lubricated and strong, she gives
birth with ease and without complications.

\item[5]

In the ninth month, one should have her enter the maternity ward.

On an auspicious lunar day and under favourable astrological
conditions.

In that context, the sacrificial hall for the priests, warriors,
merchants, and labourers should be situated on white, red, yellow, and
black plots of land respectively.

The entire chamber should be constructed from the wood of the Bilva,
Nyagrodha, or Tinduka trees; the couch therein should be made of the
same material, the walls and furnishings should be well-plastered and
arranged, and it should feature a door facing either the east or the
north, measuring eight cubits in length and four cubits in width, and be
equipped with protective and auspicious charms.

\item[6] \emph{vulgate only}

\item[7]

In that situation, a woman who has reached the time of childbirth, is
having pain all around her hips and back.

Urine is discharged frequently.

And phlegm [is discharged] from the opening of the vagina.

\item[8]

One should then have the woman who is about to give birth---who has been
well-anointed, whose limbs have been sprinkled with warm water, and for
whom auspicious rites have been performed---drink a barley gruel mixed
with clarified butter until she is full. Then, she should be placed on
a soft, wide bed provided with a pillow, lying on her back with her
thighs flexed, while three or four elderly women who are trustworthy
attend to her. It is said that women skilled in midwifery and having
trimmed nails should attend [to her].

\item[9]

Then, he should lubricate her vaginal canal in the direction of the hair
growth to ensure her comfort. And he should say to her: 
\begin{quote}
“O fortunate
one, you should now strain.

“But do not strain when the labour pains have not yet arrived.”
\end{quote}

Straining when the labour pains are not present \ldots

\ldots{} produces a child who is deaf, mute, has a dropped jaw, suffers
from head injuries, or is afflicted by cough, asthma, and wasting, or is
hunchbacked or deformed.

Therefore, she should not strain when the labour pains have not
arrived.

However, when the abdomen has descended, the foetal connections are
released, and the roots of the pelvis, groins, and bladder are relaxed,
she should strain slowly; then, when the foetus has reached the vaginal
opening, she should strain more forcefully to ensure a delivery free of
obstruction.

\item[10]

In that case, one should administer a counteractive after-drink and
cause the area to dry.

\item[11] \emph{vulgate only}

\item[12]

After cleansing the foetal membranes from the newborn and its mouth with
a mixture of rock salt and clarified butter, one should tie the
umbilical cord, leaving eight finger-breadths remaining.

Having anointed it with oil infused with \emph{costus}, one should tie
it (= the umbilical cord) with a thread and fasten it securely around
the infant's neck.

\item[13]

Then, after reviving the infant with cold water and performing the birth
rites, one should have the child lick a mixture of honey, clarified
butter, and Anantā powder along with Brāhmī juice and gold using the
ring finger.

The woman who has just given birth should be made to drink a dose of
either clarified butter or oil, followed by a decoction of the
Pippalyādi group of drugs.

One should administer gruel for three or five nights to a woman,
depending on whether she is strong or weak.

Thereafter, one should treat her with a gradual introduction of unctuous
foods, and she should be regularly administered oily substances.

\item[14--17] \emph{vulgate only}

\item[18]

Generally, one should bathe her with a generous amount of warm water and
she must abstain from anger, physical exertion, sexual intercourse, and
similar activities.

\item[19]

And there is a verse on this subject:

\begin{sloka}
    Any disease that arises in a woman after childbirth due to improper
conduct becomes difficult to cure or even incurable because of the
depletion of the bodily tissues.
\end{sloka}

\item[20]

Therefore, one should always administer treatment after carefully
examining the location, the time, the nature of the disease, and the
suitability of the procedure; in this way, one will not encounter any
failure.

\item[21]

If the placenta does not descend, one should cause it to fall; one
should rub her throat (= the cervix = throat of the uterus?) with a
finger wrapped in hair, or fumigate the vaginal opening with bitter
gourd, bottle gourd,~\gls{kṛtavedhana}, mustard seeds, and the slough
of a snake, mixed with pungent oil.

Or, one should smear the palms of her hands and the soles of her feet
with a paste of~\gls{lāṅgalī}~roots. One should make her drink a paste
of~\gls{kuṣṭha}~and~\gls{lāṅgalī}~roots mixed with either wine or
urine.

Or, a paste of~\gls{śāli}~roots or a paste of
the \gls{pippalī} group should be given with wine.

One should administer an~\emph{āsthāpana}~enema with the scum of wine
mixed with white mustard.

And with white mustard oil prepared with these same ingredients, one
should give her a vaginal douche or a powdered enema. And one should
remove it with the hand.

\item[22]

In a woman who has recently given birth, the vāta humour, by obstructing
the blood, produces pain in the bladder, the head of the bladder, the
abdomen, and the vagina; this condition is known as \emph{makkalla}.

In that case, one should administer powdered \gls{yavakṣāra} (alkali of
barley) with clarified butter and lukewarm water.

Or one should administer powdered salts with the supernatant scum of
wine.

Or she should eat old jaggery mixed with the \gls{trikaṭu}, the \gls{trijātaka}, and
coriander.

One should administer a decoction of the \gls{pṛthakparṇī} group of drugs
combined with \gls{devadāru}, black pepper, and \gls{vacā}.

\item[23]

Then, the infant, dressed in linen, should be laid upon a bed spread
with linen and cotton.

And one should fan him with branches of Nimba, Pilu, and Parushaka.

And one should apply a cotton pad soaked in oil to his head.

And one should fumigate him with incense that destroys evil spirits.

One should also fasten protective, demon-destroying herbs to his hands,
feet, head, and neck.

And one should scatter grains of sesame, linseed, and mustard here.

Furthermore, one should keep a fire burning in the room.

And one should observe the measures prescribed for the care of a patient
with an ulcer.

\item[24]

Then, on the tenth day, the mother and father, having performed
auspicious rites and protective ceremonies, should bestow a name
preceded by the recitation of benedictions.

They may choose a name according to their preference or one based upon
the birth star.

\item[25]

Then, one should engage a wet-nurse of the same caste.

She should be neither too thin nor too fat, in the prime of her youth,
free from disease, having pure milk, with breasts that are neither
pendulous nor too high, without deformity, free from vice, with a living
child, affectionate, a good milker, not mean, born of a good family, and
endowed with most of these qualities for the sake of the
child's health and strength.

In that regard, those who are excessively thin, excessively fat, or
diseased are excluded, as they cause diseases in the child due to the
imbalance of their humours.

A woman with upward-pointing breasts might make the child terrifying,
while one with pendulous breasts may cause death by covering the nose
and mouth.

Now, having seated the praised wet-nurse, who has bathed her head and is
wearing great garments, facing the north, one should first make the
child drink from the right breast.

One should consecrate it with this mantra.

\item[26]

May the four sacred oceans, which are the streams of milk from the
breasts, always provide strength to your child, O fortunate lady.

O blessed one, may they constantly grant power and vigour to the
infant.

\item[27]

May the child drink your milk, which has the flavour of nectar, O lady
of auspicious face; may he attain a long life, just as the gods did
after consuming the nectar of immortality.

May he obtain a long life, just as the gods, having consumed the nectar,
attained it.

\item[28--30] \emph{vulgate only}

\item[31]

Now, for the purpose of stimulating her milk production, one should
prescribe barley, wheat, \gls{śāli} rice, \gls{śaṣṭika} rice, meat, milk, wine,
\gls{sauvīraka}, oil-cake, garlic, fish, \gls{kaśeruka}, \gls{śṛṅgātaka}, \gls{vidārī}-bulbs,
\gls{śatāvarī}, lotus fibers, lotus stalks, \gls{alābu}, \gls{kala-shaka}, and so forth.

Then, one should examine her breast milk in water; if it is cool, pure,
thin, with the luster of a conch shell, and when placed in water it
becomes one with it, is not frothy, is not stringy, does not float, or
sinks, one should know that it is pure.

By that, the health of the child occurs.

\item[32]

And on this subject, there is a verse:~

Due to the nurse's heavy, fatty, and incompatible foods,
the humours (doṣas) in the body become aggravated, and consequently, the
breast milk becomes vitiated.

\item[33--36] \emph{vulgate only}

\item[37]

The body becomes afflicted with diseases due to the previously mentioned
corruption.

In those cases, one should administer the prescribed medicine---which
must be mild and not depleting---in an appropriate dose, combined with
milk and clarified butter, to the milk-fed infant. And to the
wet-nurse, it should be given alone.

For the child who eats both milk and solid food, medicine is prepared as
a paste.

And for the wet-nurse, it is as before. For the child who eats solid
food only, decoctions and other medicines should be given to the child
alone, and not to the wet-nurse.

\item[38--43] \emph{vulgate only}

\item[44]

And there is a verse on this subject:

\begin{sloka}
    In cases of anal suppuration in children, one should apply treatments
that alleviate the pitta humour.

Rasāñjana is particularly beneficial when used for both internal
ingestion and external application.
\end{sloka}a

\item[45]

For a child whose diet consists solely of milk, one should administer
clarified butter prepared with white mustard, sweet flag, milk-hedge,
Indian sarsaparilla, water hyssop, costus, and rock salt. For a child
who consumes both milk and solid food, clarified butter prepared with
licorice, sweet flag, and the three myrobalans should be given.

For a child who eats solid food, clarified butter prepared with the two
groups of five roots, milk, valerian, Himalayan cedar, false black
pepper, and grapes etc. should be given; by this, the
child's health, strength, intellect, and lifespan are
enhanced.

\item[46]

Furthermore, one should not seize a child by the limbs.

One should not scold him, nor wake him suddenly, for fear of causing a
fright.

One should not snatch him away suddenly or toss him up for fear of a
wind-disorder, nor force him to sit up for fear of hunchbackedness.

One should always accommodate him with things that are pleasant and
beneficial; even if he is not yet eating solid food, he will grow,
always endowed with high spirits and free from disease.

\item[47]

And there is a verse on this subject:~

One should not leave a child in an impure place, nor in the open air,
nor on uneven ground.

Not in heat, wind, rains, or in dust, smoke, or water(s).

\item[48]

Since milk is suitable for their constitution, one should give children
either goat's milk or cow's milk in
appropriate portions when breast milk is unavailable.

\item[49]

And after six months, one should feed the child food that is light and
wholesome.

\item[50]

There is a constant obstruction.

One should perform the ritual protection.

Children must be protected from the dangers of calamities and the
afflictions caused by possessing spirits.

\item[51]

Now, the boy becomes agitated, trembles, cries, and loses
consciousness.

He pricks the nurse and himself with his nails and teeth, gnashes his
teeth, moans, yawns, twitches his eyebrows, looks upward, vomits foam,
bites his lips, and possesses a harsh, broken, and piteous voice; he
stays awake at night, grows weak with withered limbs, and no longer
desires the nurse's breast milk as he did before---these
are described as the general symptoms of one afflicted by a demonic
possession.

The general symptoms of a child afflicted by a demonic seizure are
described as follows: the infant cries with a harsh, broken, or pained
voice, stays awake at night, has weak limbs with hair standing on end,
and, as previously mentioned, shows no desire for the
nurse's breast milk. "We shall explain [this] in
detail later (i.e., further on)."

\item[52]

Recognizing him to be capable, one should have him instructed in
knowledge according to his own specific status.

\item[53]

Now, for a man of twenty-five years, one should bring a wife of sixteen
years, thinking: "I shall obtain the duties of a friend, wealth, desire,
and progeny.\footnote{See discussion of this passage by \citet{wuja-2025e}.}

\item[54]

And there is a verse on this subject:
\begin{sloka}
    If a man who has not yet reached twenty-five years of age begets a
foetus with a woman who is less than sixteen years old, that foetus
perishes while still in the womb.
\end{sloka}

\item[55]

Or, having been born, he might not live for long, or he may have weak
senses.

\item[56] \emph{vulgate only}

\item[57.1]

In that situation, the foetus will fall due to the previously mentioned
causes.

Pains in the uterus, pelvis, groin, and bladder occur.

And there is the appearance of blood.

One should treat these women with sprinklings, immersions, and
ointments, as well as with the drinking of milk.

If the milk flows only very slightly, one should cause her to drink milk
processed with the \gls{utpalādi} group of drugs.

When the foetus is slipping down, there is pain in the sides with wind,
abdominal distension, and retention of urine, and the foetus moves from
one place to another.

In that case, cooling and soothing treatments should be applied.

In case of pain, one should cause her to drink milk processed with
\gls{mahāsahā}, \gls{madhuka}, \gls{svadaṃṣṭrā}, and \gls{kaṇṭakārikā}, mixed with honey and
sugar.

In case of urinary obstruction, milk processed with \gls{darbha} grass and the
like should be given.

In case of abdominal distension, a decoction of \gls{hiṅgu} and \gls{sauvarcala} is
used; if blood flows from the abdomen, one should cause her to drink a
powder of \gls{kaṇṭakārikā}, a lump of clay, \gls{gairika}, and \gls{sarjarasāñjana}, or a
paste of the bark and sprouts of the Nyagrodhādi group, as available.

Or a paste of the Utpalādi group, or a paste of \gls{kaśeru}, \gls{śṛṅgāṭaka}, and
\gls{śālūka} with boiled milk; or \gls{śāli} rice flour mixed with sugar and honey
along with a decoction of Udumbara fruit and aquatic tubers; or one
should place a piece of cloth soaked in the expressed juice of the
Nyagrodhādi group into the vagina.

Now, in case of pain without visible bleeding, one should cause her to
drink milk processed with \gls{madhuka}, \gls{devadāru}, \gls{mañjiṣṭhā}, and \gls{payasyā}; or
similarly, milk processed with \gls{aśmantaka}, \gls{śatāvarī}, and \gls{payasyā}, or
processed with the \gls{vidārigandhādi} group.

Or [milk] may be prepared with the two types of \gls{bṛhatī}, \gls{utpala},
\gls{sārivā}, \gls{payasyā}, and \gls{madhuka}.

By this [treatment], the pain is relieved and the foetus is also
nourished.


\subsection{Month-by-month development}

\item[58]

From this point forward, we shall explain the month-by-month sequels
(development of the infant).

\item[59]

Licorice, rice seeds, milk-hedge, and Himalayan cedar; mountain ebony,
black sesame, red creeper, and asparagus\ldots.

\item[60]

\gls{vṛkṣādanī}, \gls{payasyā}, \gls{latā}, \gls{utpalaśārivā}, \gls{anantā}, \gls{śārivā}, \gls{rāsnā}, lotus,
and \gls{madhuka}.

anantā, śārivā, rāsnā, lotus, and also madhuka.

\item[61]

The two types of \gls{bṛhatī}, \gls{kāśmarya}, the buds and bark of the milk-trees,
ghee, \gls{pṛthakparṇī}, \gls{balā}, \gls{śigru}, \gls{śvadaṃṣṭrā}, and \gls{madhuyaṣṭikā}.



\item[62]

Water chestnut, lotus stalk, grapes, water-nut, licorice, and sugar---O
child, these are the seven formulations, each completed in a
half-verse.

In cases of threatened miscarriage, they should be administered with
milk in their respective order.

\item[63]

The physician should administer the roots of wood apple, bael, Indian
nightshade, pointed gourd, sugarcane, and yellow-berried nightshade,
boiled in milk, during the eighth month.

In the eighth month, the physician should have the woman drink the roots
prepared with milk.

\item[64]

In the ninth month, she should drink a decoction of licorice, Indian
sarsaparilla, milk-hedge, and white sarsaparilla.

However, milk that has been boiled and then cooled is highly
recommended.

\item[65]

Alternatively, ginger, licorice, and Himalayan cedar taken with milk are
beneficial.

In this manner, the foetus is nourished and the acute pain is relieved.

\item[57.2]

Once the foetus is established, one should feed her with milk processed
with the tender sprouts of the Ficus racemosa.

When the time has passed, one should treat her with gruels made of
\gls{uddāluka} and other grains, prepared with digestive stimulants and
excluding salt, for as many days as the months of pregnancy; if there is
pain in the abdomen or foetus, she should drink an \gls{ariṣṭa} liquor
combined with old jaggery and digestive stimulants.

In a woman who indulges in wind, sun, or physical exertion, the foetus
is dissolved because the channels are constricted by the aggravated
wind; remaining there beyond the proper time, it perishes.

One should treat her with mild measures such as oleation and the like in
proper sequence.

One should provide a gruel prepared with the juice of the \gls{utkrośa} bird
and mixed with a small amount of fat.

When the foetus remains beyond the allotted time, one should
specifically pound grain with a mortar and pestle.

Or she should resort to uneven vehicles and seats.

The foetus, afflicted by wind, withers and does not fill the
mother's womb.

It moves slowly; one should treat that condition with nourishing milks
and meat juices.

Semen and blood, obstructed by wind and without an indwelling soul,
cause the abdomen to distend; if this occasionally subsides
spontaneously, they say it was “carried away by Naigameṣa.”

Sometimes, when that very foetus remains concealed or is being filled,
they call it “Nāgodara;” in that case too, the treatment is like that for
a dissolved foetus.

The child of a woman who gives birth again after a cessation of
childbearing for more than six years becomes short-lived.

\item[66] \emph{vulgate only}

\item[67]

Now, if a dangerous illness arises in a pregnant woman, one should
induce emesis using food prepared with sweet and sour flavours.

One should also encourage the downward movement of the bowels and
prescribe mild measures.

This applies to both food and drink.

Furthermore, one should eat food of mild potency and employ therapeutic
procedures that are not harmful to the foetus.

\item[68]

One should prepare a mixture of gold,~\gls{varṇacū}~(perhaps a specific
herb or powder), costus, honey, sweet flag, and clarified butter; or
alternatively, use~\gls{matsyākṣaka},~\gls{śaṅkhapuṣpī}, ghee, and
gold.

\item[69]

The flowers of the~\gls{arka}~plant, honey, clarified butter, and
powdered gold; gold dust, lapis lazuli, white~\gls{dūrvā}~grass, ghee,
and honey.

Gold dust, lapis lazuli, white~\gls{dūrvā}~grass, ghee, and honey.

\item[70]

Four auspicious medicinal pastes have been described, each completed
within half a verse, which promote the physical form, memory, strength,
and health of children.

These are auspicious for children, producing physical beauty,
intelligence, strength, and health.

\item[70.1]

Thus ends the tenth chapter on the anatomy of the body, titled "The
Prophecy of the Pregnant Woman."

The reflection on the elements, the purification of the semen, as well
as the transmission of the embryo, the prophecy of the foetus, and the
exposition of the body.

\item[70.2]

The classification of the vital points and veins, as well as the
piercing of those veins, the detailed description of the
arteries---these ten are completed within the pregnant woman.

Thus ends the section on anatomy.

\end{tt}
\end{translation}